\documentclass[crop,tikz]{standalone}% 'crop' is the default for v1.0, before it was 'preview'
%\usetikzlibrary{...}% tikz package already loaded by 'tikz' option

\usepackage{amsmath}

\usetikzlibrary{math, calc, decorations.pathreplacing}
\usetikzlibrary{shadings}
\usetikzlibrary{patterns}

\begin{document}

\tikzmath{
	% Ellipse properties.
	\a = 4;  % Semi-major axis.
	\b = 1.9;  % Semi-minor axis.
	\e2 = 1 - \b * \b / (\a * \a);  % Eccentricity squared.
	%\e = sqrt(\e2);
%	\a2 = \a * \a;  % This kills the whole compilation.
	% The point.
	\px = 1.0;
	\py = 0.7;
	\psiDeg = atan2(\py, \px);
	\psiRad = rad(\psiDeg);
    \rhoJO = sqrt(\px * \px + \py * \py) / \a;
    \rhoInv = 1 / \rhoJO;
    \cospsi = cos(\psiDeg);
    \sinpsi = sin(\psiDeg);
% Vermeille Cartesian -> geodetic, complete version
\eVal = sqrt(\e2);
\eFourVal = \e2*\e2;
\aSqVal = \a*\a;
\bSqVal = \b*\b;
\invASqVal = 1/\aSqVal;
\bSqOverAFourVal = \bSqVal/(\aSqVal*\aSqVal);
\invSixVal = 1/6;
\eTwoBOverACubedVal = \e2*\b/(\aSqVal*\a);
\eTwoHalfVal = \e2/2;
\bSqOverASqVal = \bSqVal/\aSqVal;
\bOverEVal = \b/\eVal;
\eTwoOverBVal = \e2/\b;
\bOverAVal = \b/\a;
\tSqVal = \px*\px;
\tVal = sqrt(\tSqVal);
\pVal = \tSqVal*\invASqVal;
\qVal = \py*\py*\bSqOverAFourVal;
\rVal = (\pVal + \qVal - \eFourVal)*\invSixVal;
\rTimesEightCubedVal = 8*\rVal*\rVal*\rVal;
\evVal = \rTimesEightCubedVal + \eFourVal*\pVal*\qVal;
if \evVal > 0 then {
  \sVal = \eTwoBOverACubedVal*abs(\py)*\tVal;
  \sqrtEvVal = sqrt(\evVal);
  \cOneVal = \sqrtEvVal + \sVal;
  \cTwoVal = \sqrtEvVal - \sVal;
  \uVal = \rVal + 0.5*(\cOneVal*\cOneVal)^(1/3) + 0.5*(\cTwoVal*\cTwoVal)^(1/3);
  \vVal = sqrt(\uVal*\uVal + \eFourVal*\qVal);
  \wVal = \eTwoHalfVal*(\uVal + \vVal - \qVal)/\vVal;
  \kVal = (\uVal + \vVal)/(sqrt(\wVal*\wVal + \uVal + \vVal) + \wVal);
  \DVal = \kVal*\tVal/(\kVal + \e2);
  \distVal = sqrt(\DVal*\DVal + \py*\py);
  \height = (\kVal - \bSqOverASqVal)*\distVal/\kVal;
  \latDeg = 2*atan(\py/(\distVal + \DVal));
} else {
  if \qVal > 0 then {
    \sVal = \eTwoBOverACubedVal*abs(\py)*\tVal;
    \upDeg = (2/3)*atan(\sVal/(sqrt(-\evVal) + sqrt(-\rTimesEightCubedVal)));
    \uVal = -4*\rVal*sin(\upDeg)*cos(30 + \upDeg);
    \vVal = sqrt(\uVal*\uVal + \eFourVal*\qVal);
    \wVal = \eTwoHalfVal*(\uVal + \vVal - \qVal)/\vVal;
    \kVal = (\uVal + \vVal)/(sqrt(\wVal*\wVal + \uVal + \vVal) + \wVal);
    \DVal = \kVal*\tVal/(\kVal + \e2);
    \distVal = sqrt(\DVal*\DVal + \py*\py);
    \height = (\kVal - \bSqOverASqVal)*\distVal/\kVal;
    \latDeg = 2*atan(\py/(\distVal + \DVal));
  } else {
    \height = -\bOverEVal*sqrt(\e2 - \pVal);
    \latDeg = 2*atan(sqrt(\eFourVal - \pVal)/(-\eTwoOverBVal*\height + \bOverAVal*sqrt(\pVal)));
  };
};
\latRad = rad(\latDeg);
    \phiDeg = \latDeg; % deg(\phiRad);
    \sinphi = sin(\phiDeg);
    \cosphi = cos(\phiDeg);
    \alt = \height;
    %
    coordinate \footpoint;
    \footpoint = ({\px - \alt * \cosphi}, {\py - \alt *\sinphi});
    coordinate \axispoint;
    \axispoint = ({\px - \py / tan(\phiDeg)}, 0);
    %
    \ae = \a * sqrt(\e2);
    \alphaRad = ln(1/sqrt(\e2) + sqrt(1/(\e2) - 1)); % arcosh(1/e)
    \alphaDeg = deg(\alphaRad);
    \Ce = \a * \e2 * sin(\alphaDeg) * cos(\alphaDeg)
      / sqrt(1 - \e2 * cos(\alphaDeg)^2);
    \xCap = \a+1; % desired x cutoff
    %\phiMaxBoundDeg = acos(\Ce / (\a * sqrt(\e2))) - \alphaDeg;
    \phiMaxXDeg = atan((cos(\alphaDeg) - \Ce/\xCap)/sin(\alphaDeg));
    %
    % Intersection angle between evolute and circle rho = a e
\betaVal = (\b/\a)^(2/3);
\eTwoThirdVal = (sqrt(\e2))^(2/3);
%
\psiLo = 0;
\psiHi = 90;
%
for \iter in {1,...,40} {
  \psiMid = 0.5*(\psiLo + \psiHi);
  \fMid = cos(\psiMid)^(2/3) + \betaVal*sin(\psiMid)^(2/3) - \eTwoThirdVal;
  if \fMid > 0 then {
    \psiLo = \psiMid;
  } else {
    \psiHi = \psiMid;
  };
};
%
\psiEvolAeDeg = 0.5*(\psiLo + \psiHi);
% Secant-evolute intersection
\psiLo = 0;
\psiHi = \psiEvolAeDeg;
for \iter in {1,...,50} {
  \psiMid = 0.5*(\psiLo + \psiHi);
  \rhoSecMid = \Ce/cos(\alphaDeg + \psiMid);
  \rhoEvolMid = \a*\e2 / ((cos(\psiMid)^(2/3) + (1-\e2)^(1/3)*sin(\psiMid)^(2/3))^(3/2));
  \fMid = \rhoSecMid - \rhoEvolMid;
  if \fMid > 0 then {
    \psiHi = \psiMid;
  } else {
    \psiLo = \psiMid;
  };
};
\psiIntDeg = 0.5*(\psiLo + \psiHi);
}
    
%\tikzmath{
%	\rho = 1;
%	\sin_psi = \rho * sin(60);
%	\cos_psi = \rho * cos(60);
%}

\begin{tikzpicture}[line width=0.25mm,
brc/.style = {decorate, thin, decoration={brace,raise=0.75mm, amplitude=1mm}},
shorten <>/.style = {shorten >=#1pt, shorten <=#1pt}
]
\coordinate (O) at (0,0);
\coordinate (p) at (\px cm,\py cm);

\clip (-0.5cm,-0.5cm) rectangle (\a cm+2cm, \b cm + 1.75cm);

% Fill region of convergence.
\fill[black!10]
%  ({0}:{\Ce/cos(\alphaDeg)})                         % secant meets x-axis
%  plot[domain=0:\psiIntDeg,samples=200,variable=\t]
%    ({\t}:{\Ce/cos(\alphaDeg+\t)})                   % secant bound
%  --
%  plot[domain=\psiIntDeg:\psiEvolAeDeg,samples=200,variable=\psi]
  plot[domain=0:\psiEvolAeDeg,samples=200,variable=\psi]
    ({\psi}:{\a*\e2/((cos(\psi)^(2/3) + (1-\e2)^(1/3)*sin(\psi)^(2/3))^(3/2))})
  --
  ({\psiEvolAeDeg}:{\ae}) -|
%  plot[domain=\psiEvolAeDeg:90,samples=200,variable=\t]
%    ({\t}:{\ae})                                     % circle rho = ae
%  --
  (0,0) -- cycle;
  \begin{scope}
\clip
%  ({0}:{\Ce/cos(\alphaDeg)})                         % secant meets x-axis
%  plot[domain=0:\psiIntDeg,samples=200,variable=\t]
%    ({\t}:{\Ce/cos(\alphaDeg+\t)})                   % secant bound
%  --
%  plot[domain=\psiIntDeg:\psiEvolAeDeg,samples=200,variable=\psi]
  plot[domain=0:\psiEvolAeDeg,samples=200,variable=\psi]
    ({\psi}:{\a*\e2/((cos(\psi)^(2/3) + (1-\e2)^(1/3)*sin(\psi)^(2/3))^(3/2))})
  --
  ({\psiEvolAeDeg}:{\ae}) -|
%  plot[domain=\psiEvolAeDeg:90,samples=200,variable=\t]
%    ({\t}:{\ae})                                     % circle rho = ae
%  --
  (0,0) -- cycle;
  \shade[shading=axis, top color=white, bottom color=black!10] ($({\psiEvolAeDeg}:{\ae}) + (0.4 cm, 0)$) rectangle (0, \ae cm - 0.4 cm);
  \end{scope}

% Shaded regions in quadrant 2-4.
% 2nd quadrant
\begin{scope}[xscale=-1]
\clip
%  ({0}:{\Ce/cos(\alphaDeg)})                         % secant meets x-axis
%  plot[domain=0:\psiIntDeg,samples=200,variable=\t]
%    ({\t}:{\Ce/cos(\alphaDeg+\t)})                   % secant bound
%  --
%  plot[domain=\psiIntDeg:\psiEvolAeDeg,samples=200,variable=\psi]
  plot[domain=0:\psiEvolAeDeg,samples=200,variable=\psi]
    ({\psi}:{\a*\e2/((cos(\psi)^(2/3) + (1-\e2)^(1/3)*sin(\psi)^(2/3))^(3/2))})
  --
  ({\psiEvolAeDeg}:{\ae}) -|
%  plot[domain=\psiEvolAeDeg:90,samples=200,variable=\t]
%    ({\t}:{\ae})                                     % circle rho = ae
%  --
  (0,0) -- cycle;
  % Shade from y-axis and left.
  \shade[shading=axis, left color=white, right color=black!10] (0,0) rectangle (0.4cm, \ae cm);
  % Shade up.
  \shade[shading=axis, shading angle=135, top color=white, bottom color=black!10, blend mode=lighten] ($({\psiEvolAeDeg}:{\ae}) + (-0.3 cm, 0)$) rectangle (0, \ae cm - 0.4 cm);
\end{scope}
% 3rd quadrant
\begin{scope}[xscale=-1, yscale=-1]
\clip
%  ({0}:{\Ce/cos(\alphaDeg)})                         % secant meets x-axis
%  plot[domain=0:\psiIntDeg,samples=200,variable=\t]
%    ({\t}:{\Ce/cos(\alphaDeg+\t)})                   % secant bound
%  --
%  plot[domain=\psiIntDeg:\psiEvolAeDeg,samples=200,variable=\psi]
  plot[domain=0:\psiEvolAeDeg,samples=200,variable=\psi]
    ({\psi}:{\a*\e2/((cos(\psi)^(2/3) + (1-\e2)^(1/3)*sin(\psi)^(2/3))^(3/2))})
  --
  ({\psiEvolAeDeg}:{\ae}) -|
%  plot[domain=\psiEvolAeDeg:90,samples=200,variable=\t]
%    ({\t}:{\ae})                                     % circle rho = ae
%  --
  (0,0) -- cycle;
\shade[inner color=black!10, outer color=white] (0,0) circle[radius=0.4cm];
\end{scope}

% 4th quadrant
\begin{scope}[yscale=-1]
\clip
%  ({0}:{\Ce/cos(\alphaDeg)})                         % secant meets x-axis
%  plot[domain=0:\psiIntDeg,samples=200,variable=\t]
%    ({\t}:{\Ce/cos(\alphaDeg+\t)})                   % secant bound
%  --
%  plot[domain=\psiIntDeg:\psiEvolAeDeg,samples=200,variable=\psi]
  plot[domain=0:\psiEvolAeDeg,samples=200,variable=\psi]
    ({\psi}:{\a*\e2/((cos(\psi)^(2/3) + (1-\e2)^(1/3)*sin(\psi)^(2/3))^(3/2))})
  --
  ({\psiEvolAeDeg}:{\ae}) -|
%  plot[domain=\psiEvolAeDeg:90,samples=200,variable=\t]
%    ({\t}:{\ae})                                     % circle rho = ae
%  --
  (0,0) -- cycle;
\shade[shading=axis, bottom color=white, top color=black!10] (0,0) rectangle (\ae cm, 0.4cm);
\end{scope}
%
\node[rotate=90] at (80:\b cm + 0.175cm) {\tiny{region of convergence}};



% Coordinate systems.
\draw[->, line width=0.2mm] (-\a cm -1cm,0) -- (\a cm+1cm,0) node[anchor=north west,xshift=-1mm] {$x$};
\draw[->, line width=0.2mm] (0,-\b cm - 0.7cm) -- (0,\b cm +0.7cm) node[anchor=south east,yshift=-1mm] {$y$};
% Ellipse.
\draw (O) ellipse ({\a cm} and {\b cm});
\draw[brc] (\a cm - 0.4 mm,0) -- (0.4mm, 0) node[midway, below=1.5mm] {$a$};
\draw[brc] (0, 0.4mm) -- (0, \b cm - 0.4mm) node[midway, left=1.5mm] {$b$};
%\node at (-\a*0.5 cm + 0.5cm,-\b*0.5 cm+0.1cm) {$e^2 = 1 - b^2/a^2$};
\node (ellipse_text) at (15:\a cm + 1.1cm) {\small $\tfrac{x^2}{a^2} + \tfrac{y^2}{b^2}=1$};
\draw[->] ($(ellipse_text.south west) + (0.1cm, 0.0cm)$) -- +(-145:0.35cm);
% Point, foot point and axis point.
\node[circle, fill=black, minimum size=1mm, inner sep=0] at (p) {};
\node[anchor=west] at (p) {$(u,v)$};
\node[circle, fill=black, minimum size=1mm, inner sep=0] at (O) {};
\node[circle, fill=black, minimum size=1mm, inner sep=0] at (\axispoint) {};
\node[circle, fill=black, minimum size=1mm, inner sep=0] at (\footpoint) {};
\draw (p) -- (\footpoint) -- (\axispoint);
\draw [brc] (p) -- ($(\footpoint) + (\phiDeg:0.3mm)$) node[midway, yshift=\cosphi*3.5mm, xshift=-\sinphi*3.5mm] {$h$};
% Lines and annotations.
\draw[brc] (O) -- (p) node[midway, yshift=\cospsi*3.3mm, xshift=-\sinpsi*3.3mm] {$\rho$};
\draw (O) -- (p);
\draw[thin] ($(O) + (0.4cm, 0)$) arc (0:\psiDeg:0.4cm);
\node at ($(O) + (\psiDeg/2:0.7cm)$) {$\psi$};
\draw[thin] ($(\axispoint) + (0.4cm, 0)$) arc (0:\phiDeg:0.4cm);
\node at ($(\axispoint) + (\phiDeg/2:0.6cm)$) {$\phi$};
\draw ($(\footpoint) - (\phiDeg:2mm)$) -- ++(\phiDeg-90:2mm) -- ++(\phiDeg:2mm);

% Evolute of the ellipse
%\draw[thick]
%  plot[domain=0:90,samples=200,variable=\t]
%    ({\t}:{\a*\e2/( (cos(\t)^2)^(1/3) + (1-\e2)^(1/3)*(sin(\t)^2)^(1/3) )^(3/2)});
%\node (evolute) at (48:\b cm +1.2cm) {\scriptsize{evolute}};
\node (evolute) at (\a*0.91 cm+0.5mm,\b*1.3 cm) {\small $\left(\frac{a x}{a^2-b^2}\right)^{2/3}\!+\,\left(\frac{b y}{a^2-b^2}\right)^{2/3}\!=1$};
\draw[->] ($(evolute.west)+(0.1cm,-0.1cm)$) -- +(-170:0.4cm);
\draw[thick, dotted]
  plot[domain=0:90,samples=200,variable=\t]
    ({\a*\e2*cos(\t)^3},{\a*\e2/sqrt(1-\e2)*sin(\t)^3});
\draw[thick, dotted]
  plot[domain=0:90,samples=200,variable=\t]
    ({\a*\e2*cos(\t)^3},-{\a*\e2/sqrt(1-\e2)*sin(\t)^3});
\draw[thick, dotted]
  plot[domain=0:90,samples=200,variable=\t]
    (-{\a*\e2*cos(\t)^3},{\a*\e2/sqrt(1-\e2)*sin(\t)^3});
\draw[thick, dotted]
  plot[domain=0:90,samples=200,variable=\t]
    (-{\a*\e2*cos(\t)^3},-{\a*\e2/sqrt(1-\e2)*sin(\t)^3});
    
% Circle with radius ae.
%\node (ae) at (50:\a cm+0.2cm) {$\rho=ae$};
%\draw[->] ($(ae) + (-0.15cm, -0.15cm)$) -- +(-0.25cm,-0.25cm);
%\draw[thick, dotted] (O) circle[radius=\ae cm];

% Secant bound.
%\node (secant) at (26:\a cm+1.5cm) {$\rho=\frac{C(e)}{\cos(\alpha+\psi)}$};
%\draw[->] ($(secant) + (-160: 0.7cm)$) -- +(-160:0.5cm);
%\draw[thick,dotted]
%  plot[domain=0:\phiMaxXDeg,samples=200,variable=\t]
%    ({\t}:{\Ce/cos(\alphaDeg+\t)});
%\draw[thick,dotted]
%  plot[domain=0:\phiMaxXDeg,samples=200,variable=\t]
%    ({-\t}:{\Ce/cos(\alphaDeg+\t)});
%\draw[thick,dotted]
%  plot[domain=0:\phiMaxXDeg,samples=200,variable=\t]
%    ({180-\t}:{\Ce/cos(\alphaDeg+\t)});
%\draw[thick,dotted]
%  plot[domain=0:\phiMaxXDeg,samples=200,variable=\t]
%    ({\t-180}:{\Ce/cos(\alphaDeg+\t)});

\end{tikzpicture}

\end{document}